\paragrafos{Μεσοκάθετος ευθυγράμμου τμήματος}{Β.2.3, «Μεσοκάθετος ευθυγράμμου τμήματος»}

\begin{anagke}
Στον κήπο υπάρχουν δύο μεγάλα δέντρα, το $\mathrm{Α}$ και το $\mathrm{Β}$, και η παρέα θέλει να τοποθετήσει ένα παγκάκι ώστε να είναι σε ακριβώς ίση απόσταση από τα δύο, για να έχουν όλοι τον ίδιο ίσκιο. Ο Βασίλης προτείνει το μέσο του τμήματος $\mathrm{ΑΒ}$. «Αυτό ισαπέχει», λέει η Πολυξένη, «αλλά είναι μόνο ένα σημείο. Υπάρχουν και άλλα;».

Η Αγγελική παίρνει σχοινί μήκους $5$ μέτρων, καρφώνει το ένα άκρο στο $\mathrm{Α}$ και το άλλο στο $\mathrm{Β}$ και τεντώνει από το μέσο ένα δεύτερο σχοινί του ίδιου μήκους προς κάθε μεριά. Όπου ενώνονται, βρίσκουν σημεία που ισαπέχουν από τα δύο δέντρα. Τα σημεία ενώνονται σε μια ευθεία κάθετη στο $\mathrm{ΑΒ}$, που περνά από το μέσο του. «Αυτή είναι η μεσοκάθετος», λέει η κυρία Αλεξάνδρα.

Η παρέα ανακαλύπτει ότι όλα τα σημεία της μεσοκαθέτου ισαπέχουν από τα άκρα του τμήματος, και ότι κάθε σημείο που ισαπέχει είναι πάνω της. Θα δουν και πώς κατασκευάζεται με τους δύο κλασικούς κανόνες του Ευκλείδη: τον χάρακα χωρίς υποδιαιρέσεις και τον διαβήτη.
\end{anagke}

\begin{lexi}{μεσοκάθετος, ισαπέχω}{μεσοκάθετος, ισαπέχω}
Η \lexiword{μεσοκάθετος} είναι η «κάθετη στο μέσο»: κάθετη προς το τμήμα και μάλιστα στο μέσο του. \lexiword{Ισαπέχω} σημαίνει «απέχω ίσα» (ίση απόσταση) από δύο ή περισσότερα σημεία.
\end{lexi}

\begin{metafrasi}
Η κάθετη στο μέσο του $\mathrm{ΑΒ}$. & η μεσοκάθετος του $\mathrm{ΑΒ}$\\
Ίση απόσταση από τα $\mathrm{Α}$ και $\mathrm{Β}$. & $\mathrm{ΜΑ}=\mathrm{ΜΒ}$\\
Ο άξονας συμμετρίας του τμήματος. & η μεσοκάθετος\\
Η κατασκευή με διαβήτη. & δύο ίσοι κύκλοι, κέντρα $\mathrm{Α}$ και $\mathrm{Β}$\\
\end{metafrasi}

\tmima{Η θεωρία}

\begin{orismos}{μεσοκάθετος}
\textbf{Μεσοκάθετος} ενός ευθυγράμμου τμήματος $\mathrm{ΑΒ}$ είναι η ευθεία που είναι \textbf{κάθετη} στο $\mathrm{ΑΒ}$ και διέρχεται από το \textbf{μέσο} του.
\end{orismos}

\begin{idiotita}{μεσοκάθετος και ισαπόσταση}
(α) Κάθε σημείο της μεσοκαθέτου του $\mathrm{ΑΒ}$ ισαπέχει από τα $\mathrm{Α}$ και $\mathrm{Β}$.\\
(β) Κάθε σημείο που ισαπέχει από τα $\mathrm{Α}$ και $\mathrm{Β}$ ανήκει στη μεσοκάθετο του $\mathrm{ΑΒ}$.\\
(γ) Η μεσοκάθετος του $\mathrm{ΑΒ}$ είναι άξονας συμμετρίας του τμήματος.
\end{idiotita}

Η (γ) λέει ότι το $\mathrm{Α}$ και το $\mathrm{Β}$ είναι συμμετρικά ως προς τη μεσοκάθετο (το $\mathrm{ΑΒ}$ είναι κάθετο στον άξονα και ο άξονας περνά από το μέσο του). Η (α) προκύπτει έτσι: αν $\mathrm{Μ}$ είναι σημείο του άξονα, το συμμετρικό του είναι το ίδιο, άρα η συμμετρία φέρνει το τμήμα $\mathrm{ΜΑ}$ πάνω στο $\mathrm{ΜΒ}$ και $\mathrm{ΜΑ}=\mathrm{ΜΒ}$.

\tmimasmall{Κατασκευή με κανόνα και διαβήτη}

Ο Ευκλείδης στήριζε τις κατασκευές του σε τρεις κανόνες («αιτήματα»): (1) από δύο σημεία διέρχεται μία ευθεία, (2) ένα τμήμα προεκτείνεται απεριόριστα, (3) ο κύκλος ορίζεται από κέντρο και ακτίνα. Έτσι οι αρχαίοι Έλληνες κατασκεύαζαν σχήματα με ένα χάρακα χωρίς υποδιαιρέσεις («κανόνα») και ένα διαβήτη. Για να φέρουμε τη μεσοκάθετο:

\begin{enumerate}[leftmargin=1.6em,itemsep=0.3ex,topsep=2pt]
\item Με κέντρο το $\mathrm{Α}$ και ακτίνα μεγαλύτερη από το μισό του $\mathrm{ΑΒ}$ γράφω κύκλο.
\item Με κέντρο το $\mathrm{Β}$ και την ίδια ακτίνα γράφω δεύτερο κύκλο. Οι δύο κύκλοι τέμνονται σε δύο σημεία $\mathrm{Κ}$ και $\Lambda$.
\item Η ευθεία $\mathrm{ΚΛ}$ είναι η μεσοκάθετος του $\mathrm{ΑΒ}$.
\end{enumerate}

\begin{center}
\begin{tikzpicture}[scale=0.9]
\draw (0,0) circle (2.4); \draw (3.2,0) circle (2.4);
\fill (0,0) circle (2pt) node[below left] {$\mathrm{Α}$}; \fill (3.2,0) circle (2pt) node[below right] {$\mathrm{Β}$};
\draw[thick] (0,0)--(3.2,0);
\coordinate (K) at (1.6,1.789); \coordinate (L) at (1.6,-1.789);
\fill (K) circle (2pt) node[above] {$\mathrm{Κ}$}; \fill (L) circle (2pt) node[below] {$\Lambda$};
\draw[gblue,thick,{Stealth}-{Stealth}] (1.6,2.3)--(1.6,-2.3);
\fill (1.6,0) circle (2pt) node[above right] {$\mathrm{Μ}$};
\end{tikzpicture}
\end{center}

Το σημείο $\mathrm{Μ}$ όπου η $\mathrm{ΚΛ}$ τέμνει το $\mathrm{ΑΒ}$ είναι το μέσο του.

\begin{prosoxi}
\begin{itemize}[leftmargin=1.2em,itemsep=0.3ex,topsep=2pt]
\item Η μεσοκάθετος είναι ευθεία, όχι τμήμα: συνεχίζεται και πέρα από το μέσο.
\item Δεν αρκεί να περνά από το μέσο: πρέπει να είναι και κάθετη. Μια λοξή ευθεία από το μέσο δεν είναι μεσοκάθετος.
\item Στην κατασκευή, και οι δύο κύκλοι πρέπει να έχουν την ίδια ακτίνα, μεγαλύτερη από το μισό του $\mathrm{ΑΒ}$ (αλλιώς δεν τέμνονται).
\item Το μέσο βρίσκεται με $\mathrm{ΑΜ}=\mathrm{ΜΒ}$, αλλά η μεσοκάθετος είναι η κάθετη από αυτό στο $\mathrm{ΑΒ}$.
\end{itemize}
\end{prosoxi}

\tmima{Λυμένα παραδείγματα}

\paradeigma{Ισαπέχω από τα δύο δέντρα}
\tixero Το παγκάκι $\mathrm{Μ}$ βρίσκεται πάνω στη μεσοκάθετο του $\mathrm{ΑΒ}$ και $\mathrm{ΜΑ}=5$ μ.\par
\tithelo Η απόσταση $\mathrm{ΜΒ}$.\par
\possindeo Κάθε σημείο της μεσοκαθέτου ισαπέχει από τα άκρα.\par
\apantisi $\mathrm{ΜΒ}=5$ μ.\par
\elegxos Αν ήταν $\mathrm{ΜΒ}\neq5$, το $\mathrm{Μ}$ δεν θα ήταν στη μεσοκάθετο.

\paradeigma{Κατασκευή μεσοκαθέτου}
\tixero Ένα τμήμα $\mathrm{ΑΒ}=6$ cm.\par
\tithelo Να φέρω τη μεσοκάθετο και το μέσο του.\par
\possindeo Γράφω δύο κύκλους ακτίνας $4$ cm (μεγαλύτερης από το $3$) με κέντρα $\mathrm{Α}$ και $\mathrm{Β}$. Οι δύο κύκλοι τέμνονται στα $\mathrm{Κ}$ και $\Lambda$. Η ευθεία $\mathrm{ΚΛ}$ τέμνει το $\mathrm{ΑΒ}$ στο μέσο $\mathrm{Μ}$.\par
\apantisi $\mathrm{ΑΜ}=\mathrm{ΜΒ}=3$ cm και $\mathrm{ΚΛ}\perp\mathrm{ΑΒ}$.\par
\elegxos Μετρώ: $\mathrm{ΑΜ}=3$ cm και η γωνία στο $\mathrm{Μ}$ είναι $90^\circ$.

\paradeigma{Σημείο που ισαπέχει}
\tixero Σημείο $\mathrm{Π}$ με $\mathrm{ΠΑ}=\mathrm{ΠΒ}$.\par
\tithelo Πού βρίσκεται το $\mathrm{Π}$.\par
\possindeo Αφού ισαπέχει από τα άκρα, βρίσκεται πάνω στη μεσοκάθετο του $\mathrm{ΑΒ}$.\par
\apantisi Πάνω στη μεσοκάθετο του $\mathrm{ΑΒ}$.\par
\elegxos Το μέσο του $\mathrm{ΑΒ}$ είναι το ένα τέτοιο σημείο, αλλά υπάρχουν άπειρα.

\paradeigma{Τρία δέντρα}
\tixero Τρία δέντρα $\mathrm{Α}$, $\mathrm{Β}$, $\Gamma$ που δεν είναι συνευθειακά.\par
\tithelo Πώς βρίσκω σημείο που ισαπέχει από τα τρία.\par
\possindeo Η μεσοκάθετος του $\mathrm{ΑΒ}$ περιέχει τα σημεία που ισαπέχουν από τα $\mathrm{Α}$, $\mathrm{Β}$. Η μεσοκάθετος του $\mathrm{ΒΓ}$ τα σημεία που ισαπέχουν από τα $\mathrm{Β}$, $\Gamma$. Το σημείο τομής τους $\mathrm{Ο}$ ισαπέχει και από τα τρία.\par
\apantisi Το σημείο τομής δύο μεσοκαθέτων.\par
\elegxos Τότε $\mathrm{ΟΑ}=\mathrm{ΟΒ}=\mathrm{ΟΓ}$ και ο κύκλος με κέντρο $\mathrm{Ο}$ περνά από τα τρία δέντρα.

\tmima{Ασκήσεις}

\zoni{Το κατάλαβα}{μεσοκάθετος και ισαπόσταση}

\ex Σχεδίασε τμήμα $\mathrm{ΑΒ}=6$ cm και φέρε τη μεσοκάθετο με κανόνα και διαβήτη. Μέτρησε το $\mathrm{ΑΜ}$ και την $\widehat{\mathrm{ΑΜΚ}}$.

\ex Πάνω στη μεσοκάθετο του $\mathrm{ΑΒ}$ πάρε τρία διαφορετικά σημεία και μέτρησε τις αποστάσεις τους από τα $\mathrm{Α}$ και $\mathrm{Β}$.

\ex Αν το $\mathrm{Μ}$ ανήκει στη μεσοκάθετο του $\mathrm{ΑΒ}$ και $\mathrm{ΜΑ}=5$, πόσο είναι το $\mathrm{ΜΒ}$;

\ex Αν $\mathrm{ΠΑ}=\mathrm{ΠΒ}$, πού βρίσκεται το $\mathrm{Π}$;

\ex Σχεδίασε τρίγωνο και φέρε τις μεσοκαθέτους των δύο πλευρών του. Σημείωσε το σημείο τομής τους και μέτρησε τις αποστάσεις του από τις τρεις κορυφές.

\ex Σωστό ή λάθος; α) Η μεσοκάθετος περνά από το μέσο του τμήματος. \quad β) Κάθε ευθεία που περνά από το μέσο είναι μεσοκάθετος. \quad γ) Η μεσοκάθετος είναι άξονας συμμετρίας του τμήματος. \quad δ) Κάθε σημείο εκτός της μεσοκαθέτου ισαπέχει από τα $\mathrm{Α}$, $\mathrm{Β}$.

\zoni{Το χρησιμοποιώ}{παγκάκια, πηγάδια, δρόμοι}

\ex \textbf{Το παγκάκι.} Τα δέντρα $\mathrm{Α}$ και $\mathrm{Β}$ απέχουν $8$ μ. Βρες δύο θέσεις για το παγκάκι ώστε να απέχει εξίσου από τα δύο δέντρα και ακριβώς $5$ μ. από το καθένα.

\ex \textbf{Το πηγάδι.} Δύο χωριά $\mathrm{Α}$ και $\mathrm{Β}$ θέλουν ένα πηγάδι που να ισαπέχει από αυτά και να είναι πάνω στο ευθύγραμμο δρόμο που περνά από το μέσο τους (κάθετος στο $\mathrm{ΑΒ}$). Πού θα το βάλουν;

\ex \textbf{Οι τρεις φίλοι.} Τρεις φίλοι μένουν σε διαφορετικά σημεία και θέλουν να συναντιούνται σε σημείο που να ισαπέχει από τα σπίτια τους. Πώς θα το βρουν;

\ex \textbf{Ο κύκλος.} Βρες το κέντρο ενός κύκλου με μία χορδή του $\mathrm{ΑΒ}$ και μία δεύτερη $\mathrm{ΒΓ}$ φέρνοντας μεσοκαθέτους.

\ex \textbf{Το φανάρι.} Ένα φανάρι πρέπει να τοποθετηθεί ώστε να ισαπέχει από δύο σχολεία $\mathrm{Α}$ και $\mathrm{Β}$ και να είναι στον ευθύγραμμο δρόμο $\varepsilon$ που τέμνει το $\mathrm{ΑΒ}$. Πώς βρίσκεται η θέση;

\zoni{Το σκέφτομαι}{η αρχαία κατασκευή}

\ex Γιατί στην κατασκευή της μεσοκαθέτου οι δύο κύκλοι πρέπει να έχουν ακτίνα μεγαλύτερη από το μισό του $\mathrm{ΑΒ}$;

\ex Γιατί η μεσοκάθετος είναι άξονας συμμετρίας του τμήματος; Τι είναι το συμμετρικό του $\mathrm{Α}$;

\ex Αν δύο σημεία ισαπέχουν από τα $\mathrm{Α}$ και $\mathrm{Β}$, πώς είναι η ευθεία που τα ενώνει ως προς το $\mathrm{ΑΒ}$;

\ex Πόσα σημεία έχουν απόσταση $5$ cm και από τα δύο άκρα ενός τμήματος $\mathrm{ΑΒ}=6$ cm;

\zoni{Για το σπίτι}{κατασκευές}

\ex Φέρε με κανόνα και διαβήτη τη μεσοκάθετο ενός τμήματος και έλεγξέ την δίπλωνοντας το χαρτί ώστε το $\mathrm{Α}$ να πέσει πάνω στο $\mathrm{Β}$.

\ex Βρες το κέντρο ενός κυκλικού αντικειμένου, π.χ. ενός δίσκου, χρησιμοποιώντας δύο χορδές.
